\documentclass[11pt,a4paper]{article}
\usepackage{fontspec,amsmath,amssymb,amsthm,microtype}
\setmainfont{lmroman10-regular.otf}[BoldFont=lmroman10-bold.otf,ItalicFont=lmroman10-italic.otf,BoldItalicFont=lmroman10-bolditalic.otf]
\setmonofont{lmmono10-regular.otf}[Scale=MatchLowercase]
\usepackage[margin=25mm]{geometry}
\usepackage[unicode,hidelinks]{hyperref}
\setlength{\emergencystretch}{3em}
\setlength{\parskip}{0.4em}
\setlength{\parindent}{0pt}
\allowdisplaybreaks[2]
\raggedbottom
\title{\texorpdfstring{The true reference velocity is independent of its following Brownian interval}{The true reference velocity is independent of its following Brownian interval}}
\author{LangevinLean mathematical exposition}
\date{}
\begin{document}\maketitle
\noindent Source module: \texttt{Kinetic\allowbreak{}Langevin/\allowbreak{}Canonical\allowbreak{}Past\allowbreak{}Future.\allowbreak{}lean}. The accompanying HTML pairs this mathematical exposition with the exact Lean statements and records the scope of the formalization.


This module proves the stochastic source condition needed in the repaired second-order KLMC argument. The actual reference trajectory is constructed from its initial phase and past Brownian input. Its terminal velocity is independent of the genuine future OU convolution. Stationarity then identifies their actual joint law as standard Gaussian velocity times the exact thermal Gaussian innovation.


\section{\texorpdfstring{Make dependence on a finite innovation prefix explicit}{Make dependence on a finite innovation prefix explicit}}


A recursion at index n reads only noise entries with index less than n. Induction proves this for the general random recursion and for the deterministic Verlet–OU splitting. These identities permit arbitrary continuation of a finite past noise grid without changing its terminal state.



{\raggedright\noindent Lean declarations: \texttt{noise\allowbreak{}Driven\allowbreak{}Chain\_\allowbreak{}eq\_\allowbreak{}of\_\allowbreak{}prefix}, \texttt{kinetic\allowbreak{}Split\allowbreak{}Noise\allowbreak{}Chain\_\allowbreak{}eq\_\allowbreak{}of\_\allowbreak{}prefix}.\par}

\section{\texorpdfstring{Identify actual source representatives}{Identify actual source representatives}}


The coupling extension has independent Gaussian velocity and a canonical Brownian coordinate. The square-root fluctuation-dissipation factor is already built into its continuous linear noise lift. The representative theorem proves that each lifted segment is the original Brownian random vector, evaluated at that coordinate and multiplied by the exact factor.



{\raggedright\noindent Lean declarations: \texttt{past\allowbreak{}Future\_\allowbreak{}coupling\allowbreak{}Probability}, \texttt{input\allowbreak{}Lift\_\allowbreak{}coe}.\par}

\section{\texorpdfstring{Prove independence for genuine splitting approximants}{Prove independence for genuine splitting approximants}}


Divide the past horizon into n+1 intervals. Continue its innovation stream by repeating the last past interval, so every deterministic kernel stays entirely before the future interval. The proved Brownian support theorem makes this whole stream independent of the future response. Adjoin the independent initial positions and velocity, and reassociate the actual product source. The finite pointwise splitting recursion is a measurable function of precisely those past inputs. Its L\ensuremath{{}^{2}} representative is the original vector recursion. Prefix dependence identifies its endpoint with the actual continuous-input splitting chain, proving independence of each true approximant.



{\raggedright\noindent Lean declarations: \texttt{split\allowbreak{}Chain\_\allowbreak{}velocity\_\allowbreak{}future\_\allowbreak{}independent}.\par}

\section{\texorpdfstring{Pass to the actual nonlinear kinetic reference}{Pass to the actual nonlinear kinetic reference}}


The matching-input splitting theorem converges in L\ensuremath{{}^{2}} to the actual continuous kinetic integral solution on any finite horizon. Projecting to velocity preserves this convergence. Independence from the fixed future convolution is closed under L\ensuremath{{}^{2}} convergence, so the limiting reference velocity is independent of that same actual future random vector. The future marginal itself is exactly the translated OU thermal Gaussian law, by the proved Brownian convolution distribution and measure-preserving source projection.



{\raggedright\noindent Lean declarations: \texttt{global\allowbreak{}Reference\allowbreak{}Flow\_\allowbreak{}velocity\_\allowbreak{}future\_\allowbreak{}independent}, \texttt{future\allowbreak{}Input\_\allowbreak{}has\allowbreak{}Law}.\par}

\section{\texorpdfstring{Export the exact Gaussian pair laws}{Export the exact Gaussian pair laws}}


The actual reference phase remains Gaussian-times-Gibbs distributed by the proved stationarity theorem. Its velocity marginal is therefore standard Gaussian. Combining this marginal, the exact future thermal Gaussian marginal and the newly proved independence gives their full product joint law. The stitched reference chain equals the same continuous solution at grid points, yielding the identical joint-law statement for every numerical interval.



{\raggedright\noindent Lean declarations: \texttt{global\allowbreak{}Reference\allowbreak{}Flow\_\allowbreak{}velocity\_\allowbreak{}future\_\allowbreak{}joint\allowbreak{}Law}, \texttt{reference\allowbreak{}Chain\_\allowbreak{}velocity\_\allowbreak{}future\_\allowbreak{}joint\allowbreak{}Law}.\par}

\end{document}
